\section{Modèle local et accès à MongoDB}\label{sec:modele}
MongoDB fonctionne comme un serveur séparé dans Docker Compose ; l'application
y accède par une URI réseau. Le déploiement local utilise MongoDB 8.0.16 sur
\code{127.0.0.1:27018}, avec un volume persistant. L'ingestion et les contrôles
Parquet sont en Python. L'API est en Spring Boot 3.5.7, Java 21, avec Spring Data
MongoDB et \code{MongoTemplate}. JPA n'est pas utilisé : les requêtes sont des
filtres et pipelines BSON explicites. Les contrôleurs valident les entrées, les
services appliquent les règles et les repositories interrogent la base. Postman
sert d'interface de démonstration \cite{sujet,code}. Les accès attendus sont des filtres par zone et par
intervalle de dates suivis de regroupements : les index de la section~\ref{sec:mesures}
servent ces filtres. MongoDB offre aussi le partitionnement natif, qui laisse à l'API
une seule URI en déploiement réparti (section~\ref{sec:reparti}).

La collection \code{trips} conserve les colonnes TLC dans \code{data}, la provenance
\code{source.file_id}/\code{source.row_number}, la durée \code{duration_s} et les
flags. \code{zones} associe un identifiant entier aux noms et boroughs : les noms
ne sont pas des clés uniques. \code{imports} conserve le SHA-256 du fichier,
ses métadonnées et son état \code{loading}, \code{failed} ou \code{complete}.
Les services ne sélectionnent que les imports complets. Aucun agrégat matérialisé
n'est nécessaire aux services décrits ici.

\paragraph{Identité et reprise.}
L'identité d'un trajet combine l'heure de départ, une empreinte de toutes les
colonnes canonisées et le rang parmi les lignes identiques du fichier. Le code
utilise une empreinte Pandas de 64 bits, \emph{distincte} du SHA-256 de provenance.
Deux copies dans un fichier ont des rangs différents ; une ligne déjà importée
retrouve son identité et l'index unique \code{_id} empêche sa réinsertion.
L'import utilise des insertions par lots et traite les collisions de clé
comme des lignes déjà présentes.
Cette règle ne résout pas toute ambiguïté : deux enregistrements identiques issus
de fichiers indépendants peuvent être confondus, une correction de contenu
produit une nouvelle identité, et une collision d'empreinte reste possible.
En cas de recouvrement, le document garde la provenance de la première insertion,
pas toutes ses rééditions. La section~\ref{sec:f6} teste cette règle sur de vraies lignes.

\paragraph{Dates et durées.}
Les heures locales TLC sont stockées en chaînes ISO sans fuseau. Les conversions
UTC ne sont que temporaires pour calculer la durée : attribuer New York aux
deux heures, convertir en UTC, puis soustraire. Une arrivée inexistante ou ambiguë
produit \code{null}, pas une durée locale de secours. La ligne 942\,762 de mars
part le 9 mars à 23:38:43 et arrive le 10 à 03:40:14 : la durée réelle est
10\,891 s, contre 14\,491 s par soustraction locale. Elle porte
\code{DST_SPANNING}. La ligne 841\,170 arrive à 02:09:54 le 10 mars, heure qui
n'existe pas : elle porte \code{DST_UNRESOLVED} et est exclue de journey.

\paragraph{Calculs côté serveur et traçabilité.}
\code{$match} filtre les documents ; \code{$project} prépare les champs utiles
sans modifier la base ; \code{$group} regroupe et compte. \code{$sum:1} compte
les candidats. Pour les retenus, \code{$cond} vaut 1 si la ligne est éligible et
0 sinon, puis \code{$sum} additionne ces valeurs. La même sélection produit
ainsi les effectifs bruts et retenus. Le booléen est calculé par
l'intersection entre les flags du trajet et ceux du profil, grâce à
\code{$setIntersection}, \code{$size} et \code{$eq}. Une ligne cumulant plusieurs
motifs n'est donc exclue qu'une fois.

Les groupes incluent le fichier source pour justifier les comptes. \code{GET /api/trips} compte et liste, avec pagination,
les lignes sous-jacentes selon zones, dates, créneau, profil et statut, et renvoie
le fichier et la position d'origine.  La présence de fichiers complets ne prouve pas la
couverture de chaque heure ; les services ne fournissent pas non plus un
instantané transactionnel entre la lecture des imports et l'agrégation.
